\section{Concepts and Scope}\label{sec:background}
\figwide{Fig2_survey_structure.pdf}{The structure of this survey. Green leaves locate representative studies. Branches navigate detection signals, evaluation targets, and interventions, rather than define mutually exclusive classes. Tree design adapted from \citet{liu-etal-2025-logical-design}; content is specific to this survey.}{fig:structure}
\subsection{Reference frames}
\emph{World factuality} concerns agreement with externally established facts; \emph{source faithfulness} concerns support by supplied evidence. Neither entails the other: a response can repeat a false source faithfully or add a true but unlicensed detail \citep{cao-etal-2022-hallucinated,adlakha-etal-2024-evaluating}. BEGIN evaluates attribution in grounded dialogue under this distinction \citep{dziri-etal-2022-evaluating}.

Source-based labels should distinguish \emph{supported}, \emph{contradicted}, and \emph{not established}, retaining the last separately from world falsity. World verification additionally needs a specified evidence source, date, and unresolved-claim policy. These are operational choices, not a universal definition. HalluLens revisits definitions and benchmark construction \citep{bang-etal-2025-hallulens}; comparisons must map study-specific labels before comparing scores.

\subsection{Error relations and development}
Entity, attribute, and relation errors are orthogonal to reference frames: a wrong date can contradict either a document or a world fact. Annotation units also differ. ANAH provides sentence-level analysis, while HaluEval~2.0 extends empirical factuality analysis across domains \citep{ji-etal-2024-anah,li-etal-2024-dawn}.

Research moved from distinct task-specific correctness targets in 2022 toward atomic verification and black-box self-checking in 2023 \citep{min-etal-2023-factscore,manakul-etal-2023-selfcheckgpt}, followed by contextual monitoring and factuality alignment. Recent studies scrutinize judge reliability, information coverage, and temporal validity \citep{godbole-jia-2025-verify,samarinas-etal-2025-beyond,jiang-etal-2026-benchmarks}. This trajectory does not imply that newer methods supersede earlier ones.

We separate three axes: reference frame, detection evidence, and intervention point. A benchmark evaluates a detector; a detector can also guide decoding. Potential failure points include unfamiliar supervision, knowledge access, retrieval, and evidence use. Fine-tuning and long-context studies motivate these possibilities without proving a complete causal decomposition \citep{gekhman-etal-2024-fine,liu-etal-2024-lost}.
